\section{Hexagonal 3-webs, Blaschke curvature, singularities}
A planar 3-web ${\cal W}_3$ in a planar domain is a superposition of 3 foliations $\mathcal{F}_i$, which may be given by integral curves of three ODEs
$
\sigma _1=0$, $\sigma _2=0$, $ \sigma _3=0$,
where $\sigma _i$ are differential one-forms. At non-singular points, where the kernels of these forms are pairwise transverse, we normalize the forms so that
$\sigma_1+\sigma_2+\sigma_3=0.$ The {\it connection form} of the web ${\cal W}_3$ is a one-form $\gamma$ determined by the conditions
${\rm d}\sigma_i+\gamma \wedge \sigma_i=0$, $i=1,2,3$.
The connection form depends on the normalization of the forms $\sigma_i$, the {\it Blaschke curvature} $d\gamma$ does not.
\begin{Definition}
A 3-web is hexagonal if for any non-singular point there are a neighbourhood and a local diffeomorphism sending the web leaves of this neighbourhood in 3 families of parallel line segments.
\end{Definition}
Topologically, hexagonality means the following incidence property that has given its name to the notion: for any point $m$, each sufficiently small curvilinear triangle with the vertex $m$ and sides formed by the web leaves, may be completed to the curvilinear hexagon, whose sides are web leaves and whose ``large'' diagonals are the web leaves meeting at $m$ (see the gallery of pictures illustrating hexagonal webs in the next section). Computationally, hexagonality amounts to vanishing of the Blaschke curvature \cite{BB-38}.

Up to a suitable affine transformation, the forms $\sigma_i$ may be normalized as follows:
\[
\sigma_1=(Q - R)({\rm d}y - P{\rm d}x), \qquad \sigma_2=(R - P)({\rm d}y - Q{\rm d}x), \qquad \sigma_3=(P - Q)({\rm d}y - R{\rm d}x),
\]
where $P(x,y)$, $Q(x,y)$, $R(x,y)$ are the slopes of the tangent lines to the web leaves at $(x,y)$. Vanishing of the curvature writes as
\begin{gather}
(R-Q )[P_{xx} +(Q + R)P_{xy} +QRP_{yy}]+(P-R)[Q_{xx} +(P + R)Q_{xy}+PRQ_{yy}]\nonumber\\
\qquad{}+ (Q - P)[R_{xx} +(P + Q)R_{xy}+PQR_{yy}]\nonumber\\
\qquad{}+ \frac{(Q - R)(2P - Q - R)\bigl[P^2_x+(Q+R)P_xP_y+QRP^2_y\bigr]}{(P-Q )(P - R)}\nonumber \\
\qquad{} +
 \frac{(R-P )(2Q -P- R)\bigl[Q^2_x+(P+R)Q_xQ_y+PRQ_y^2\bigr]}{(Q - R)(Q-P)}
\nonumber \\
\qquad{}+ \frac{(P - Q)(2R-P-Q)\bigl[R_x^2+(P+Q)R_xR_y+PQR_y^2\bigr]}{(R-Q)(R - P)} \nonumber\\
\qquad{}+
\frac{(2R-P-Q)P_xQ_x }{P-Q}+
 \frac{(2P-Q-R)Q_xR_x}{Q - R}+
 \frac{(2Q-R-P)R_xP_x}{R-P} \nonumber\\
\qquad{}+ \frac{ \bigl(R^2-PQ\bigr)[P_xQ_y+P_yQ_x]}{(P - Q)} +
 \frac{\bigl(P^2 - QR\bigr)[Q_xR_y+Q_yR_x]}{Q-R} \nonumber\\
\qquad{} +
 \frac{\bigl(Q^2-PR\bigr)[R_xP_y+R_yP_x]}{R-P} + \frac{\bigl(2PQR -(P+Q)R^2\bigr)P_yQ_y}{Q-P}\nonumber\\
 \qquad{} +
\frac{\bigl(2PQR-(Q+R)P^2\bigr)Q_yR_y}{R-Q} +
\frac{\bigl(2PQR-(P+R)Q^2\bigr)R_yP_y }{P-R} =0.\label{curvature}
\end{gather}

If only one slope, say $R$, is given as an explicit functions of $x$, $y$ and $P$, $Q$ are roots of a~quadratic equation $P^2+AP+B=0$, then one finds the first derivatives of $P$, $Q$ by differentiating the Vieta relations
\begin{equation}\label{Vieta}
P+Q=-A,\qquad PQ=B,
\end{equation}
as a functions of $P$, $Q$ and the first derivatives of $A$, $B$. Differentiating these expressions, one gets also the second derivatives. Finally, excluding $P$ and $Q$ with the help of \eqref{Vieta}, one can rewrite \eqref{curvature} in terms of $A$, $B$, $R$ and their derivatives. The result is presented in Appendix~\ref{AppendixA}.

The webs considered later will inevitably have singularities: some kernels of the forms $\sigma_i$ can be not transverse or the forms can vanish at some points. We call a singular 3-web hexagonal if its Blaschke curvature vanishes identically at regular points. The simplest type of singularities of hexagonal 3-webs have the following remarkable property, first observed by Shelekhov \cite{S-05}.

\begin{Lemma}\label{sing}
Suppose that a hexagonal $3$-web, defined by three analytic direction fields $\xi_1$,~$\xi_2$,~$\xi_3$, has a singular point $p_0$ such that
\begin{enumerate}\itemsep=0pt
\item[$(1)$] all $\xi_i$ are well defined at $p_0$,
\item[$(2)$] $\xi_1$ and $\xi_2$ are transverse at $p_0$,
\item[$(3)$] $\xi_1=\xi_3$ at $p_0$,
\end{enumerate}
then either the leaves of $\xi_1$ and $\xi_3$ through $p_0$ coincide or
$\xi_1=\xi_3$ along the leaf of $\xi_2$ through $p_0$.
\end{Lemma}
\begin{proof}
The property is a consequence of separation of variables for hexagonal webs.
The second condition implies that we can rectify $\xi_1$ and $\xi_2$, i.e., choose some local coordinates $u$, $v$ so that $\xi_1=\partial_v$ and $\xi_2=\partial_u$. Then $\xi_3=-f(u,v)\partial_u+\partial_v$ with $f(u_0,v_0)=0$, where $p_0=(u_0,v_0).$
Consider the analytic functions $\varphi(u):=f(u,v_0)$.

If $\varphi(u)\equiv 0$, then $f(u,v_0)=0$ and the leaves of $\xi_1$ and $\xi_3$ are tangent along the line $v=v_0$, which is the leaf of $\xi_2$.

 If $\varphi(u)\not\equiv 0$, then $\varphi(u)=(u-u_0)^n\psi(u)$ with natural $n$ and $\psi(u_0)\ne 0$. Thus for any~$u_1$ close to $u_0$ with $u_1\ne u_0$ holds $f(u_1,v_0)=(u_1-u_0)^n\psi(u_1)\ne 0$. Now the hexagonality amounts to~${\partial_v\bigl(\frac{\partial_uf}{f}\bigr)=0}$ hence the germ $\tilde{f}$ of $f$ at $(u_1,v_0)$ factors $\tilde{f}(u,v)=a(u)b(v)$ with analytic germs $a$, $b$ at $u_1$ and $v_0$ respectively. One can choose $b(u)$ so that $b(v_0)=1$. Then $a(u)=\psi(u)=(u-u_0)^n\psi(u)$ is analytic also at $u_0$, the function of two variables $a(u)b(v)$ is analytic at $(u_0,v_0)$ and coincides with $f(u,v)$ at some neighborhood of $(u_1,v_0)$ included in the domain of $f$. Then by uniqueness $f(u,v)=a(u)b(v)$. Now observe that $a(u_0)=0$ which
implies $f(u_0,v)\equiv 0$ and the integral curves of $\xi_1$ and $\xi_3$ passing through $p_0$ coincide.
\end{proof}

\section{Projective model of M\"obius geometry}
Following Blaschke \cite{B-29}, we call the subgroup ${\rm PSO}(3,1)$ of projective transformations of $\mathbb{RP}^3$, leaving invariant the quadric
$
X^2+Y^2+Z^2-U^2=0$,
the M\"obius group.
For the reference, we present here infinitesimal generators of M\"obius group in homogeneous coordinates $[X:Y:Z:U]$ in $\mathbb{P}^3$, affine coordinates $x=\frac{X}{U}$, $y=\frac{Y}{U}$, $ z=\frac{Z}{U}$ in $\mathbb{R}^3$ and cartesian coordinates $(\bar{x},\bar{y})$ in $\mathbb{R}^2$ related to points $(x,y,z)$ on the unit sphere via stereographic projection
\smash{$
x=\frac{2\bar{x}}{1+\bar{x}^2+\bar{y}^2}$}, \smash{$ y=\frac{2\bar{y}}{1+\bar{x}^2+\bar{y}^2}$}, \smash{$z=\frac{1-\bar{x}^2-\bar{y}^2}{1+\bar{x}^2+\bar{y}^2}$}.
There are 3 rotations around the affine axes:
\begin{gather*}
R_z=Y\partial_X-X\partial_Y=y\partial_x-x\partial_y=\bar{y}\partial_{\bar{x}}-\bar{x}\partial_{\bar{y}},\\
R_x=Z\partial_Y-Y\partial_Z=z\partial_y-y\partial_z=\bar{x}\bar{y}\partial_{\bar{x}}+\frac{1}{2} \bigl(1-\bar{x}^2+\bar{y}^2\bigr)\partial_{\bar{y}},\\
R_y=X\partial_Z-Z\partial_X=x\partial_z-z\partial_x=-\frac{1}{2} \bigl(1+\bar{x}^2-\bar{y}^2\bigr)\partial_{\bar{x}}-\bar{x}\bar{y}\partial_{\bar{y}},
\end{gather*}
and 3 boosts (or ``hyperbolic rotations'')
\begin{gather*}
B_x=U\partial_X+X\partial_U=\partial_x-x(x\partial_x+y\partial_y+z\partial_z)=\frac{1}{2} \bigl(1-\bar{x}^2+\bar{y}^2\bigr)\partial_{\bar{x}}-\bar{x}\bar{y}\partial_{\bar{y}},\\
B_y=U\partial_Y+Y\partial_U=\partial_y-y(x\partial_x+y\partial_y+z\partial_z)= -\bar{x}\bar{y}\partial_{\bar{x}}+\frac{1}{2} \bigl(1+\bar{x}^2-\bar{y}^2\bigr)\partial_{\bar{y}},\\
B_z=U\partial_Z+Z\partial_U=\partial_z-z(x\partial_x+y\partial_y+z\partial_z)=-\bar{x}\partial_{\bar{x}}-\bar{y}\partial_{\bar{y}}.
\end{gather*}
The identity component of ${\rm PSO}(3,1)$ is well known to be isomorphic to the group $PSL_2(\mathbb{C})$, the isomorphism being given by the action $A(V)=AVA^*$ of $A\in SL_2(\mathbb{C})$ on the vector space of
matrices
\[
V=\left(
\begin{matrix}
X+U & Y+{\rm i}Z\\
Y-{\rm i}Z & U-X
\end{matrix}
\right)
\]
with real $X$, $Y$, $Z$, $U$. This action preserves determinant of~$V$, which is $U^2-X^2-Y^2-Z^2$.
By this isomorphism, the generators are represented by the following matrices:
\begin{gather*}
R_x=\frac{1}{2}\left(
\begin{matrix}
-{\rm i} & 0\\
0 & {\rm i}
\end{matrix}
 \right),\qquad
 R_y=\frac{1}{2}\left(
\begin{matrix}
0 & -{\rm i}\\
-{\rm i} & 0
\end{matrix}
 \right),\qquad
 R_z=\frac{1}{2}\left(
\begin{matrix}
0 & 1\\
-1 & 0
\end{matrix}
 \right),
\\
B_x={\rm i}R_x,\qquad B_y={\rm i}R_y, \qquad B_z={\rm i}R_z.
\end{gather*}

Two points $p_1=[X_1:Y_1:Z_1:U_1]$ and $p_2=[X_2:Y_2:Z_2:U_2]$ in $\mathbb{RP}^3$ determine a line with the Pl\"ucker coordinates
\begin{alignat*}{4}
& a:=X_1U_2-X_2U_1,\qquad&& b:=Y_1U_2-Y_2U_1, \qquad&& c:=Z_1U_2-Z_2U_1,&
\\
& f:=Y_1Z_2-Y_2Z_1,\qquad && g:=Z_1X_2-Z_2X_1, \qquad&& h:=X_1Y_2-X_2Y_1.&
\end{alignat*}
By direct computation, one proves the following fact.
\begin{Lemma}
All points of a line with Pl\"ucker coordinates $[a:b:c:f:g:h]$ are stable with respect to subgroup with the infinitesimal generator $aR_x+bR_y+cR_z+fB_x+gB_y+hB_z$.
\end{Lemma}
Observe that the line dual to $[a:b:c:f:g:h]$ is the one with coordinates $[-f:-g:-h:a:b:c]$, which corresponds to multiplication by ${\rm i}$ of the corresponding matrix representation of the generator.

A line in $\mathbb{RP}^3$ can be hyperbolic, elliptic or parabolic with the respect to the Darboux quadric. Considering simple representatives of these classes one sees that
\begin{itemize}\itemsep=0pt
\item[(1)] For hyperbolic line, the action of the corresponding operator on the dual line is M\"obius conjugate to the action of $R_z$ on the line $Z=U=0$ at infinity thus not having extra stable points in $\mathbb{RP}^3$.

\item[(2)] For elliptic line, the action of the corresponding operator on the dual line is M\"obius conjugate to the action of $B_z$ on the line $X=Y=0$ (which action leaves invariant two extra points $[0:0:\pm1:1]$).

\item[(3)] For parabolic line, the corresponding operator moves all points on the dual line (consider~${\partial_y=R_x+B_y}$).
\end{itemize}

In the following proposition we summarize further properties of the above correspondence.
\begin{Proposition}\label{operline}
Let $\xi =aR_x+bR_y+cR_z+fB_x+gB_y+hB_z$ be an infinitesimal operator of the M\"obius group.
\begin{enumerate}\itemsep=0pt
\item[$(1)$] The corresponding action of the one-parameter group is not loxodromic and therefore M\"obius equivalent $($i.e., conjugated$)$ either to rotation, or dilatation, or translation of the plane $(\bar{x},\bar{y})$ if and only if \begin{equation}\label{Pluck}
af+bg+ch=0.
\end{equation}
\item[$(2)$] This action is M\"obius equivalent to rotation if and only if \eqref{Pluck} and
$a^2+b^2+c^2>f^2+g^2+h^2$ are true, the line with Pl\"ucker coordinates $[a:b:c:f:g:h]$ being the set of polar points of the circular orbits.
\item[$(3)$] This action is M\"obius equivalent to dilatation if and only if~\eqref{Pluck} and
$a^2+b^2+c^2<f^2+g^2+h^2$ are true, the line with Pl\"ucker coordinates $[a:b:c:f:g:h]$ being the set of polar points of the orbits.
\item[$(4)$] The action is M\"obius equivalent to translation if and only if \eqref{Pluck} and
\begin{equation}\label{quadrcomplex}
a^2+b^2+c^2=f^2+g^2+h^2
\end{equation} are true, the line with Pl\"ucker coordinates $[a:b:c:f:g:h]$ being the set of polar points of the orbits.
\end{enumerate}
\end{Proposition}
\begin{proof}
The type of subgroup transformations is determined by the eigenvalues of the matrix representation for the generating operator $\xi$. Consider the characteristic polynomial for $\xi$
\[
\Psi (\lambda)=\lambda^2 +\frac{1}{4}\bigl( a^2 + b^2 + c^2 - f^2 - g^2 - h^2\bigr)+\frac{\rm i}{2}(af + bg +ch)
\]
and the matrices for non-loxodromic M\"obius representatives $R_z$, $B_z$ and $\partial_y=R_x+B_y$.
\end{proof}
